% theme: uniform_circular_motion_and_centripetal_force
\MajorQuestion{等速円運動と向心力}
\SetRowSpace{2mm}

\TextTall{円周率は $\pi$ とする。必要に応じて、等速円運動の速さを $v$、角速度を $\omega$、半径を $r$、周期を $T$、質量を $m$ として考えなさい。}

\TypeQuestion{質量 $0.25\,{\rm kg}$ の小球が、半径 $0.50\,{\rm m}$ の円周上を角速度 $4.0\,{\rm rad/s}$ で等速円運動している。次の問いに答えなさい。}{angular_velocity_to_motion_quantities}
\QQ{小球の速さを求めなさい。}{$2.0\,{\rm m/s}$}
\QQ{小球の向心加速度の大きさを求めなさい。}{$8.0\,{\rm m/s^2}$}
\QQ{小球にはたらく合力の大きさを求めなさい。}{$2.0\,{\rm N}$}
\QQ{小球の周期を求めなさい。}{$\dfrac{\pi}{2}\,{\rm s}$}

\TypeQuestion{質量 $0.50\,{\rm kg}$ の物体が、半径 $2.0\,{\rm m}$ の円周上を速さ $6.0\,{\rm m/s}$ で等速円運動している。次の問いに答えなさい。}{speed_to_motion_quantities}
\QQ{物体の角速度を求めなさい。}{$3.0\,{\rm rad/s}$}
\QQ{物体の向心加速度の大きさを求めなさい。}{$18\,{\rm m/s^2}$}
\QQ{物体にはたらく合力の大きさを求めなさい。}{$9.0\,{\rm N}$}
\QQ{物体の周期を求めなさい。}{$\dfrac{2\pi}{3}\,{\rm s}$}

\TypeQuestion{質量 $0.20\,{\rm kg}$ の小球が、半径 $0.25\,{\rm m}$ の円周上を周期 $0.50\,{\rm s}$ で等速円運動している。次の問いに答えなさい。}{period_to_motion_quantities}
\QQ{小球の振動数を求めなさい。}{$2.0\,{\rm Hz}$}
\QQ{小球の角速度を求めなさい。}{$4\pi\,{\rm rad/s}$}
\QQ{小球の速さを求めなさい。}{$\pi\,{\rm m/s}$}
\QQ{小球の向心加速度の大きさを求めなさい。}{$4\pi^2\,{\rm m/s^2}$}

\AnswerColumnBreak
\LessonQuestionBreak

\TextTall{重力加速度の大きさを $g$ とし、空気抵抗は無視する。必要に応じて、物体にはたらく力を図示して考えなさい。}
\SetRowSpace{0mm}

\TypeQuestion{質量 $m$ の小球を長さ $\ell$ の軽くて伸びない糸につけ、糸が鉛直方向となす角を $\theta$ に保ったまま、角速度 $\omega$ で水平面内を等速円運動させた。糸の張力を $S$ とする。次の問いに答えなさい。}{conical_pendulum}
\FigProblemBlock{34mm}{fig/snipe_01_conical_pendulum}{
\QQ{小球が描く円の半径 $r$ を、$\ell$ と $\theta$ を用いて表しなさい。}{$r=\ell\sin\theta$}
\QQ{鉛直方向の力のつり合いを表す式を書きなさい。}{$S\cos\theta=mg$}
\QQ{水平方向について、小球の運動方程式を書きなさい。}{$S\sin\theta=m\omega^2r$}
\QQ{小球の角速度 $\omega$ を、$g$、$\ell$、$\theta$ を用いて表しなさい。}{$\omega=\sqrt{\dfrac{g}{\ell\cos\theta}}$}
}{
}

\TypeQuestion{}{motion_on_smooth_conical_surface}
\FigProblemBlock{34mm}{fig/snipe_01_conical_surface}{\Text{中心軸が鉛直な、なめらかな円錐面の内側に質量 $m$ の小球を置く。円錐面の母線は水平面と角 $\theta$ をなし、小球は中心軸から距離 $r$ の水平な円周上を速さ $v$ で等速円運動している。円錐面から小球が受ける垂直抗力の大きさを $N$ とする。次の問いに答えなさい。}}{
\QQ{鉛直方向の力のつり合いを表す式を書きなさい。}{$N\cos\theta=mg$}
\QQ{水平方向について、小球の運動方程式を書きなさい。}{$N\sin\theta=m\dfrac{v^2}{r}$}
\QQ{垂直抗力の大きさ $N$ を、$m$、$g$、$\theta$ を用いて表しなさい。}{$N=\dfrac{mg}{\cos\theta}$}
\QQ{小球の速さ $v$ を、$r$、$g$、$\theta$ を用いて表しなさい。}{$v=\sqrt{rg\tan\theta}$}
}

\TypeQuestion{}{object_on_horizontal_turntable}
\FigProblemBlock{34mm}{fig/snipe_01_turntable}{\Text{水平な回転台の上に質量 $m$ の小物体を置き、回転軸から距離 $r$ の位置で、角速度 $\omega$ で回転させた。回転台と物体の間の静止摩擦係数を $\mu$ とし、物体は回転台に対して静止している。垂直抗力の大きさを $N$、物体にはたらく静止摩擦力の大きさを $f$ とする。次の問いに答えなさい。}}{
\QQ{垂直抗力の大きさ $N$ を、$m$ と $g$ を用いて表しなさい。}{$N=mg$}
\QQ{静止摩擦力の大きさ $f$ を、$m$、$r$、$\omega$ を用いて表しなさい。}{$f=mr\omega^2$}
\QQ{角速度の最大値 $\omega_{\max}$ を、$\mu$、$g$、$r$ を用いて表しなさい。}{$\omega_{\max}=\sqrt{\dfrac{\mu g}{r}}$}
\QQ{速さの最大値 $v_{\max}$ を、$\mu$、$g$、$r$ を用いて表しなさい。}{$v_{\max}=\sqrt{\mu gr}$}
}
